% from ega0-19-fr.tex
\oldpage[0\textsubscript{IV-1}]{99}

cet homomorphisme se factorise en
$r:C\xrightarrow{r'}G\to G/\mathfrak R$, où $r'$ est un
$A$-homomorphisme continu, si bien que $G$ est ainsi muni d'une structure de
$C$-algèbre topologique (discrète). D'autre part, par composition avec
l'homomorphisme canonique
\[
  \mathfrak J/\mathfrak J^2\longrightarrow
  \mathfrak J_\beta/\mathfrak J_\beta^2\longrightarrow F_{\beta,n},
\]
on déduit de (19.5.6.9) un $C$-homomorphisme continu
$g:\mathfrak J/\mathfrak J^2\to G/\mathfrak R$. Comme
$\mathfrak J/\mathfrak J^2$ est un $C$-module formellement projectif, $g$ se
factorise en
\[
  \mathfrak J/\mathfrak J^2\xrightarrow{h}G\longrightarrow G/\mathfrak R,
\]
où $h$ est un $C$-homomorphisme continu. Comme $G$ est discret, il existe un
indice $\gamma\geq\beta$ tel que l'image par $h$ de
$((\mathfrak b_\gamma\cap\mathfrak J)+\mathfrak J^2)/\mathfrak J^2$ soit
nulle, ainsi que l'image par $r'$ de
$(\mathfrak b_\gamma+\mathfrak J)/\mathfrak J$, de sorte que $h$ se factorise
en
\[
  \mathfrak J/\mathfrak J^2\longrightarrow
  \mathfrak J_\gamma/\mathfrak J_\gamma^2
  \xrightarrow{h'}G,
\]
où $h'$ est un $C_\gamma$-homomorphisme. Par prolongement, on déduit donc de
$h'$ un homomorphisme de $C_\gamma$-algèbres
\[
  w:\mathbf S_{C_\gamma}^*(\mathfrak J_\gamma/\mathfrak J_\gamma^2)
  \longrightarrow G
\]
et par construction, tout élément de degré $m$ a pour image par $w$ un élément
de $\mathfrak R$, donc tout élément de degré $n=2m$ a une image nulle, puisque
$\mathfrak R^2=0$~; autrement dit, $w$ se factorise en
\[
  \mathbf S_{C_\gamma}^*(\mathfrak J_\gamma/\mathfrak J_\gamma^2)
  \longrightarrow F_{\gamma,n}\xrightarrow{w'_\gamma}G.
\]
Appliquons maintenant à
\[
  v_{\alpha\gamma}:F_{\gamma,n}
  \xrightarrow{q_{\beta\gamma}}F_{\beta,n}
  \xrightarrow{v_{\alpha\beta}}E_{\alpha,n}
\]
le lemme (19.5.6.7), dont les hypothèses sont vérifiées~; il existe donc un
$\delta\geq\gamma$ et un homomorphisme
$u_{\gamma\delta}:E_{\delta,n}\to F_{\gamma,n}$ tel que le composé
$v_{\alpha\gamma}\circ u_{\gamma\delta}$ soit l'homomorphisme de transition
$p_{\alpha\delta}:E_{\delta,n}\to E_{\alpha,n}$. On obtient donc finalement un
diagramme commutatif de $A$-homomorphismes continus
\[
\xymatrix@C=3em@R=2.5em{
  && E_{\alpha,n}\ar[r]^{f_{\alpha,n}} & G/\mathfrak R \\
  && F_{\beta,n}\ar[u]_{v_{\alpha\beta}} & \\
  B\ar[r]_{p_\delta} & E_{\delta,n}\ar[uur]^{p_{\alpha\delta}}
    \ar[r]_{u_{\gamma\delta}} &
  F_{\gamma,n}\ar[u]_{q_{\beta\gamma}}\ar[r]^{w'_\gamma} & G\ar[uu]
}
\]
et le composé $f':B\to G$ des homomorphismes de la ligne inférieure est donc
tel que $f$ se factorise en $B\xrightarrow{f'}G\to G/\mathfrak R$, ce qui
prouve que $B$ est une $A$-algèbre formellement lisse.

La démonstration du th.~(19.5.3) est ainsi achevée.
\end{env}

\begin{corollary}[19.5.7]
\label{0.19.5.7-fr}
Soient $A$ un anneau topologique, $B$ une $A$-algèbre topologique,
$(\mathfrak b_\lambda)$ un système fondamental d'idéaux ouverts dans $B$,
$\mathfrak J$ un idéal de $B$, $C=B/\mathfrak J$ la $A$-algèbre topologique
quotient. On pose $C_\lambda=B/(\mathfrak b_\lambda+\mathfrak J)$. On suppose
que~: $1^\circ$ pour tout $n$, la topologie induite sur $\mathfrak J^n$ par
celle de $B$ est aussi la topologie du $B$-module $\mathfrak J^n$ déduite de
la topologie de $B$ (19.0.2) (cette condition sera remplie en particulier si
$B$ est noethérien et sa topologie préadique ($\mathbf 0_{\mathrm I}$,
7.3.2))~; $2^\circ$ $C$ est une $A$-algèbre formellement lisse. Sous ces
conditions~:
\begin{enumerate}
  \item[\rm(i)] Si $B$ est une $A$-algèbre formellement lisse, alors, pour tout
  $\lambda$, $(\mathfrak J/\mathfrak J^2)\otimes_C C_\lambda$ est un
  $C_\lambda$-module projectif.
  \item[\rm(ii)] Si $B$ est une $A$-algèbre formellement lisse et un anneau
  préadmissible, alors pour tout $\lambda$, l'homomorphisme canonique
  \begin{equation}
  \label{0.19.5.7.1-fr}
    \varphi_\lambda=\varphi\otimes 1_{C_\lambda}:
    \mathbf S_{C_\lambda}^*((\mathfrak J/\mathfrak J^2)\otimes_C C_\lambda)
    \longrightarrow\gr_{\mathfrak J}^*(B)\otimes_C C_\lambda
    \tag{19.5.7.1}
  \end{equation}
  est bijectif.

  